\insertImage{4d3ea27d-bc3c-4d12-999c-9c5b794ad904.png}{Aquarelle représentant diverses parties prenantes en cercle autour d'une personne inspirante représentant une entreprise libérée.}

\chapter{Les impacts sur les parties prenantes externes}

Une entreprise libérée adore parler d'elle-même — c'est presque un symptôme. Ses équipes, ses rituels, sa culture, sa transformation. Mais aucune entreprise ne vit en vase clos : il y a des clients qui appellent, des fournisseurs qui livrent, des partenaires qui signent, et un territoire qui regarde tout ça d'un œil parfois perplexe.

Ce chapitre sort donc des murs. Que voit-on de l'extérieur quand une entreprise se libère ? Et qu'est-ce que ça change, concrètement, pour ceux qui traitent avec elle ?

\section{Relation avec les clients dans une entreprise libérée}

Pour un client, une entreprise libérée se reconnaît à un détail qui change tout : la personne en face de vous peut dire oui.

Dans une organisation classique, le premier interlocuteur est rarement le décideur — il transmet, il escalade, il revient vers vous. Dans une entreprise libérée, celui qui traite votre problème a le pouvoir de le résoudre. Chez Chrono Flex, c'est le cœur du réacteur : le technicien qui intervient sur votre chantier décide sur place, sans appeler un bureau\footnote{Gérard, A. (2017). Le patron qui ne voulait plus être chef. Flammarion.}. Pour le client pressé — et dans le dépannage hydraulique, le client est toujours pressé —, la différence se mesure en heures de production sauvées.

Chez FAVI, la logique était poussée jusqu'à l'architecture même de l'entreprise : des mini-usines dédiées chacune à un client, avec un contact direct entre les opérateurs et « leur » client\footnote{Getz, I. (2009). Liberating Leadership: How the Initiative-Freeing Radical Organizational Form Has Been Successfully Adopted. California Management Review, 51(4), 32-58.}. L'ouvrier qui coule le laiton sait pour qui il le coule. Difficile de faire plus court comme circuit qualité.

Soyons complets : ce niveau de personnalisation a ses fragilités. Il dépend de la constance des équipes — un client habitué à l'excellence pardonne mal les jours sans. Et certains clients, notamment les grandes organisations, sont déroutés par l'absence d'interlocuteur hiérarchique : « qui est le responsable ? » n'est pas toujours une question de méfiance, c'est parfois juste une case obligatoire de leur processus achats. Une entreprise libérée doit apprendre à répondre à cette question sans renier son modèle — généralement en désignant des référents, ce qui n'est pas la même chose que des chefs.

\section{Relation avec les fournisseurs et les partenaires commerciaux}

La libération déborde naturellement sur la chaîne de valeur, pour une raison de simple cohérence : on ne peut pas prôner la confiance en interne et pressurer ses fournisseurs comme au bon vieux temps.

Les entreprises libérées tendent donc vers des relations de partenariat plutôt que des rapports de force : transparence sur les contraintes, co-construction des solutions, engagement dans la durée. Ce n'est pas de l'angélisme — c'est le même pari que le modèle fait en interne, appliqué à l'externe : la confiance coûte moins cher que le contrôle, à condition d'être réciproque.

Mais la réciprocité, justement, ne se décrète pas plus dehors que dedans. Tous les fournisseurs ne veulent pas « co-construire » — certains veulent un bon de commande précis et être payés à l'heure, et c'est une position parfaitement respectable. Et tous les partenaires ne jouent pas le jeu : la transparence unilatérale, face à un acheteur resté à l'ancienne école, peut devenir un handicap de négociation. Comme pour les salariés, la règle devrait être la même : le partenariat est une offre, pas une imposition.


\section{Impact sur la communauté locale et la société}

C'est peut-être l'effet le moins commenté et le plus profond du modèle : ce qu'une entreprise libérée fait au territoire qui l'entoure.

L'exemple le plus parlant reste FAVI, et il ne s'agit pas d'une opération de mécénat. Pendant que la sous-traitance automobile délocalisait en masse, FAVI est restée à Hallencourt, village picard de quelques centaines d'habitants, en faisant le pari que des équipes responsabilisées seraient plus compétitives que des coûts salariaux réduits\footnote{Getz, I. (2009). Liberating Leadership: How the Initiative-Freeing Radical Organizational Form Has Been Successfully Adopted. California Management Review, 51(4), 32-58.}. Pour un territoire rural, une entreprise industrielle qui reste — et qui embauche —, c'est un impact social qui se passe de discours.

Il y a aussi un effet plus diffus, que je crois réel même s'il est difficile à mesurer : des gens qui apprennent au travail à décider, à se coordonner sans chef, à porter des responsabilités, ne redeviennent pas passifs en franchissant la porte de l'entreprise. Le travail est une école — la question est seulement de savoir ce qu'il enseigne. Une organisation qui enseigne l'initiative et la coopération rend à la société des citoyens un peu différents de celle qui enseigne l'obéissance et la prudence.

Gardons-nous cependant d'un raccourci : « libérée » ne veut pas dire « vertueuse ». Une entreprise peut être libérée et polluer, être libérée et vendre n'importe quoi. L'engagement sociétal et environnemental est un choix distinct, qui demande ses propres preuves. Inversement, des organisations qui ne se revendiquent pas du modèle — je pense aux coopératives, dont le pouvoir partagé est inscrit dans les statuts mêmes — produisent depuis longtemps des effets sociaux comparables sans jamais employer le mot « libération »\footnote{Rothschild, J., \& Whitt, J. A. (1986). The cooperative workplace. Cambridge University Press.}. Le modèle n'a pas le monopole du sens.

Au fond, la question que ce chapitre pose aux dirigeants tentés par la libération est celle-ci : pour qui libérez-vous ? Si la réponse s'arrête aux murs de l'entreprise, vous passez à côté de la moitié de l'histoire.
